\section{方式三：RGB-D实时检测与避障跟随（第3组）}
\subsection{方案设计}
本组采用RGB-D信息实时寻找并跟随TonyPi。程序包括颜色与深度候选检测、多帧关联、外观模板、历史轨迹目标、滚动代价地图和速度门控。目标与环境信息主要来自当前传感器和本车里程计，不要求提前建好全局地图。\cite{follow3}

现有说明将检测阈值作为示例，要求登记真实身份模板并在独立样本上评估。区域深度有效不等于识别到了TonyPi；颜色候选也不能替代身份连续性验证。源码和合成演练已提供，真实目标识别、导航与全队终点验收仍需对应组确认。

\subsection{RGB-D观测与身份关联}
感知节点将RGB与配准深度按时间同步，以颜色候选或可选的深度人体轮廓提出目标区域。对候选执行面积、有效深度比例、深度离散程度和物理尺寸筛选；颜色区域深度采用有限样本统计，不将无效深度视为远处自由空间。设候选像素$(u,v)$深度为$z$，相机内参为$f_x,f_y,c_x,c_y$，其光学坐标近似为
\begin{equation}
X=(u-c_x)z/f_x,\qquad Y=(v-c_y)z/f_y,\qquad Z=z.
\end{equation}
点云经采样及TF变换用于目标与局部环境表达；该几何过程要求深度尺度、配准和相机模型正确，不能仅凭有彩色框判断几何链可靠。

身份关联结合位置连续性和已登记外观模板，在多个候选接近时进行歧义判断。默认模板为空，只允许影子诊断；真机执行需登记TonyPi正面、背面、侧面及抱物等实际视角，并用独立负样本检查相似物体误识别。跟踪确认与重识别分别需要连续观测，不能把一次颜色匹配直接作为新身份。目标中断后不将未知位置直接连为已观察轨迹。

\subsection{本车里程计下的跟随与主动避障}
有效目标变换至本车\code{odom}，以历史路径回退点作为导航目标。与方式二相比，方式三不依赖预先保存的全局地图和AMCL，但更依赖里程计连续性、局部深度覆盖及身份稳定。滚动代价地图依据实时点云更新，未知区域保守处理；深度缺测不能被错误转为空闲清除射线。

源码为目标丢失、任务过期、TF/里程计/代价地图无效设置停止与取消导航分支，并对导航失败、里程计跳变设置需复位的故障状态。重捕获确认后重建轨迹，避免沿遮挡期间的虚构路径前进。视野约束用于判断运动后目标是否仍可能处于相机可观察范围，不能据此保证遮挡不会发生。

\subsection{速度门控、制动边界与终点}
独立门控结合命令和许可时效、观测覆盖、障碍点云以及速度制动包络限制运动。配置中的最大线速度0.12\,m/s、角速度0.4\,rad/s、回退间距1.0\,m与最小欧氏距离0.65\,m均为默认起点。制动模型需要实测减速度和反应时间，不能用导航栈加速度限制直接宣称物理停车距离。

完成逻辑同时考虑任务done、目标稳定、回退点容差、车辆足迹进入标定的终点多边形以及低速持续保持。默认\code{finish_polygon}为空，\code{hardware_calibrated}、\code{identity_calibrated}和\code{depth_coverage_calibrated}均为false，表明当前共享配置并非已完成硬件标定的比赛配置。静止判断阈值和保持2\,s为程序条件，实际到达仍需现场记录支持。

\subsection{与其他组的配合及测试安排}
本组需要与第1组一起确认TonyPi不同方向、不同抱物姿态下的外观，并结合其行走速度调整跟随参数。与队列组交流时，还要明确首车实际使用的程序和目标信息，不能只根据“点云跟踪”的名称判断两者实现相同。

我们将合成输入和只显示结果的调试模式用于检查数据链，再通过实际图像登记目标特征。完整测试还应包括静止识别、直行、转角遮挡、相似物体和动态障碍，记录目标误切换、丢失时间与停车情况。当前现场测试还不充分，这部分需要继续补充。

\subsection{小组合作}
何奇涛牵头RGB-D输入和检测，林子昂牵头目标关联与跟随，李雨翰牵头局部地图和速度保护，陈韦诺负责试验记录与整理。四人共同参与样本采集和联调，具体分工见第\ref{sec:division}节。我们需要把感知画面、深度和机器人运动放在一起检查，避免只在单一模块中反复调参。
